% Options for packages loaded elsewhere
\PassOptionsToPackage{unicode}{hyperref}
\PassOptionsToPackage{hyphens}{url}
\documentclass[
  11pt,
  letterpaper,
  twocolumn,
]{article}
\usepackage{xcolor}
\usepackage[margin=0.85in,columnsep=0.3in]{geometry}
\usepackage{amsmath,amssymb}
\setcounter{secnumdepth}{-\maxdimen} % remove section numbering
\usepackage{iftex}
\ifPDFTeX
  \usepackage[T1]{fontenc}
  \usepackage[utf8]{inputenc}
  \usepackage{textcomp} % provide euro and other symbols
\else % if luatex or xetex
  \usepackage{unicode-math} % this also loads fontspec
  \defaultfontfeatures{Scale=MatchLowercase}
  \defaultfontfeatures[\rmfamily]{Ligatures=TeX,Scale=1}
\fi
\usepackage{lmodern}
\ifPDFTeX\else
  % xetex/luatex font selection
\fi
% Use upquote if available, for straight quotes in verbatim environments
\IfFileExists{upquote.sty}{\usepackage{upquote}}{}
\IfFileExists{microtype.sty}{% use microtype if available
  \usepackage[]{microtype}
  \UseMicrotypeSet[protrusion]{basicmath} % disable protrusion for tt fonts
}{}
\makeatletter
\@ifundefined{KOMAClassName}{% if non-KOMA class
  \IfFileExists{parskip.sty}{%
    \usepackage{parskip}
  }{% else
    \setlength{\parindent}{0pt}
    \setlength{\parskip}{6pt plus 2pt minus 1pt}}
}{% if KOMA class
  \KOMAoptions{parskip=half}}
\makeatother
\usepackage{longtable,booktabs,array}
\newcounter{none} % for unnumbered tables
\usepackage{calc} % for calculating minipage widths
% Correct order of tables after \paragraph or \subparagraph
\usepackage{etoolbox}
\makeatletter
\patchcmd\longtable{\par}{\if@noskipsec\mbox{}\fi\par}{}{}
\makeatother
% Allow footnotes in longtable head/foot
\IfFileExists{footnotehyper.sty}{\usepackage{footnotehyper}}{\usepackage{footnote}}
\makesavenoteenv{longtable}
\setlength{\emergencystretch}{6em} % prevent overfull lines
\providecommand{\tightlist}{%
  \setlength{\itemsep}{0pt}\setlength{\parskip}{0pt}}
\usepackage{amsmath, amssymb, amsthm}
\usepackage{booktabs}
\usepackage{microtype}
\usepackage[hidelinks,colorlinks=true,linkcolor=blue!50!black,urlcolor=blue!50!black,citecolor=blue!50!black]{hyperref}
\hypersetup{
  pdftitle={A High-Precision Phenomenological Relation for Newton's Constant},
  pdfauthor={Oldřich Dvořák},
  pdfsubject={Phenomenological physics, gravitation, fine-structure constant},
  pdfkeywords={Newton constant, fine-structure constant, B3, SO(7), SO(8), phenomenology, mechanism gap}
}
\setlength{\emergencystretch}{6em}
\providecommand{\tightlist}{\setlength{\itemsep}{0pt}\setlength{\parskip}{0pt}}
\usepackage{bookmark}
\IfFileExists{xurl.sty}{\usepackage{xurl}}{} % add URL line breaks if available
\urlstyle{same}
\hypersetup{
  hidelinks,
  pdfcreator={LaTeX via pandoc}}

\title{A High-Precision Phenomenological Relation for Newton's Constant\\[4pt]
\large with Exact $B_3/SO(7)/SO(8)$ Algebraic Scaffold and an Explicit Mechanism Gap}
\author{Oldřich Dvořák\\[2pt]
\small Independent researcher\\
\small \texttt{oldrich@oldrich.me}\\
\small DOI: \url{https://doi.org/10.5281/zenodo.20120946}\\
\small Public proof package: \url{https://github.com/oldrich-research/gravitational-constant-relation}}
\date{May 2026}

\begin{document}

\twocolumn[%
\maketitle

\begin{center}
\begin{minipage}{0.88\textwidth}
\setlength{\parskip}{8pt plus 2pt minus 1pt}
\setlength{\parindent}{0pt}
\begin{center}\large\textbf{Abstract}\end{center}
\vspace{2pt}

\noindent
The relation
\[\boxed{\;G \;=\; \frac{4}{3}\,\frac{\hbar c}{m_e^2}\,\alpha^{21}\,\exp\!\left(-\frac{5\alpha}{2}\right)\;}\]
evaluates on CODATA 2022 inputs {[}1{]} to
\(G_{\rm pred} = 6.6743123\!\times\!10^{-11}\,\mathrm{m^3\,kg^{-1}\,s^{-2}}\), residual \(+1.84\) ppm (\(z=+0.082\sigma\)). Against the April 2026 NIST/Metrologia single-experiment result {[}10{]}, residual \(+66.3\) ppm (\(+1.16\sigma\)). Propagated input uncertainty: \(3.3\!\times\!10^{-9}\), four orders below the experimental floor.

The discrete coefficients identify exactly with \(B_3/SO(7)/SO(8)\) Lie-algebraic data: \(21\!=\!\dim(\mathrm{adj}\,SO(7))\); \(5/2\!=\!h^\vee(B_3)/2\); \(4/3\!=\!\dim(SO(8))/\dim(SO(7))\).

A bounded reproducibility scan over \(22{,}680\) candidates of class \(A(\hbar c/m_e^2)\alpha^n\exp(-k\alpha)\) places \((4/3,21,5/2)\) at rank \(1\) within \(0.5\sigma\) and \(0.1\sigma\); nearest competitor \(-3640\) ppm at \(-162\sigma\). A Julia bounded-equation-search engine enumerates \(2.02\!\times\!10^9\) raw expressions (\(2.46\!\times\!10^6\) canonical-unique) through cost shell \(C\!\le\!20\); with \(\mathrm{alpha\_family}\) pre-seeded, no competitor beats the canonical formula at budgets \(30\) and \(40\). Recovered rank-\(1\) \(\mathrm{rel\_err}=1.85\!\times\!10^{-6}\), invariant across \(C\!=\!18,19,20\). Sham-parity: \(0/256\) at \(C\!=\!18\), \(1/256\) at \(C\!=\!19,20\) (single non-\(\alpha\) TRASH expression).

A structured survey across QED, soft-photon and Sudakov resummation, instantons, functional determinants on \(S^7\) and \(B_3\), heat-kernel localization, one-loop \(\beta\)-coefficients, \(SO(7)^\pm\) gauged supergravity, induced gravity, and the Kalinski QED-vacuum kernel produces no mechanism reproducing the \(\alpha^{21}\) counting, \(\exp(-5\alpha/2)\) damping, and \(m_e\) normalization simultaneously. The relation is a high-precision phenomenological anomaly with an exact algebraic scaffold and a documented mechanism gap. Operational falsification exposure: \(\sim\!13.25\) ppm at \(5\sigma\) against any future low-uncertainty \(G\) recommendation converging on the NIST 2026 central value.

\vspace{8pt}
\noindent\small\textit{AI authorship.} This paper reports research carried out by an orchestration of AI agents under the named author's direction. The author is an independent researcher without specialist training in theoretical physics; his role was vision, target selection, and approval. See \emph{Statement of AI Authorship} for the full disclosure.

\end{minipage}
\end{center}
\vspace{1.5em}
]

\section{Headline Quantities}\label{headline-quantities}

\begin{table*}[!htb]
\centering\small
\begin{tabular}{@{}
  >{\raggedright\arraybackslash}p{(\linewidth - 2\tabcolsep) * \real{0.5000}}
  >{\raggedright\arraybackslash}p{(\linewidth - 2\tabcolsep) * \real{0.5000}}@{}}
\toprule
\begin{minipage}[b]{\linewidth}\raggedright
Quantity
\end{minipage} & \begin{minipage}[b]{\linewidth}\raggedright
Value
\end{minipage} \\
\midrule
\(G_{\rm pred}\) and residuals &
\(6.6743123 \times 10^{-11}\,\mathrm{m^3\,kg^{-1}\,s^{-2}}\); \(+1.84\)
ppm at \(+0.082\sigma\) vs CODATA 2022 {[}1{]}; \(+66.3\) ppm at
\(+1.16\sigma\) vs NIST 2026 {[}10{]} \\
Bounded prior scan (\(22{,}680\) candidates) & canonical
\((4/3, 21, 5/2)\) at rank \(1\) within both \(0.5\sigma\) and
\(0.1\sigma\); nearest non-canonical row \(-3640\) ppm at
\(-162\sigma\) \\
AV6 symbolic engine, \(C\!\le\!20\) & \(2.02\!\times\!10^9\) raw,
\(2.46\!\times\!10^6\) canonical-unique; canonical rank \(1\) at budgets
\(30, 40\); \(\mathrm{rel\_err}=1.85\!\times\!10^{-6}\) invariant across
\(C\!=\!18,19,20\) \\
Sham-parity check (\(256\) targets, budget \(40\)) & \(0/256\) at
\(C\!=\!18\); \(1/256\) at \(C\!=\!19,20\) (non-\(\alpha\) TRASH) \\
Exact algebraic identities & \(21\!=\!\dim(\mathrm{adj}\,SO(7))\);
\(5/2\!=\!h^\vee(B_3)/2\!=\!\rho\!\cdot\!e_1\);
\(4/3\!=\!\dim(SO(8))/\dim(SO(7))\) \\
Falsification threshold & \(\sim\!13.25\) ppm at \(5\sigma\) \\
\bottomrule
\end{tabular}
\end{table*}

\section{1. The Relation and Its Numerical
Status}\label{the-relation-and-its-numerical-status}

The relation evaluates to
\(G_{\rm pred} = 6.6743123 \times 10^{-11}\,\mathrm{m^3\,kg^{-1}\,s^{-2}}\)
on CODATA 2022 inputs {[}1{]}. Against the recommended value
\(G = 6.67430(15) \times 10^{-11}\) this is \(+1.84\) ppm and
\(z\!=\!+0.082\sigma\). The integer exponent required to hit the CODATA
2022 central value exactly is \(n_{\rm exact} = 21.00000037\).

Cross-benchmark status against the April 16, 2026 NIST/Metrologia
single-experiment determination \(G = 6.67387(38) \times 10^{-11}\)
{[}10{]} is \(+66.3\) ppm and \(z\!=\!+1.16\sigma\).

Propagation of CODATA input uncertainties gives a relative uncertainty
in \(G_{\rm pred}\) of \(3.3 \times 10^{-9}\), dominated \(96.5\%\) by
\(\alpha\) and \(3.5\%\) by \(m_e\). This is below the experimental
uncertainty on \(G\) by a factor of \(6.9 \times 10^3\). \(\hbar\) and
\(c\) are exact by SI definition. Recomputation across successive CODATA
releases gives \(z = +0.75\sigma\) (2014), \(+0.082\sigma\) (2018),
\(+0.082\sigma\) (2022): the 2014→2018 improvement was driven by
movement of the experimental central value toward the formula.

\section{2. Bounded-Family Reproducibility
Search}\label{bounded-family-reproducibility-search}

The search class is

\[G = A\,\frac{\hbar c}{M^2}\,\alpha^n\,\exp(-k\alpha)\]

with declared parameters fixed as follows:

\begin{itemize}
\tightlist
\item
  Mass scale \(M = m_e\) (CODATA 2022 {[}1{]}).
\item
  Exponent window \(n \in \{5,6,\dots,40\}\) (\(36\) values).
\item
  Damping window \(k \in \{0,\,1/2,\,1,\,\dots,\,10\}\) (\(21\)
  half-integer values).
\item
  Prefactor prior: the reduced-rational set generated by \(p/q\) with
  \(p\!\le\!5\), \(q\!\le\!30\), \(\gcd(p,q)\!=\!1\),
  \(p/q\!\in\![1/3,5]\). Explicit list (30 values):
  \begin{equation*}
  \begin{aligned}
  \{ &\tfrac{1}{3},\ \tfrac{5}{14},\ \tfrac{4}{11},\ \tfrac{3}{8},\ \tfrac{5}{13},\ \tfrac{2}{5}, \\
  &\tfrac{5}{12},\ \tfrac{3}{7},\ \tfrac{4}{9},\ \tfrac{5}{11},\ \tfrac{1}{2},\ \tfrac{5}{9}, \\
  &\tfrac{4}{7},\ \tfrac{3}{5},\ \tfrac{5}{8},\ \tfrac{2}{3},\ \tfrac{5}{7},\ \tfrac{3}{4}, \\
  &\tfrac{4}{5},\ \tfrac{5}{6},\ 1,\ \tfrac{5}{4},\ \tfrac{4}{3},\ \tfrac{3}{2}, \\
  &\tfrac{5}{3},\ 2,\ \tfrac{5}{2},\ 3,\ 4,\ 5 \}.
  \end{aligned}
  \end{equation*}
\item
  Total cardinality: \(30 \times 36 \times 21 = 22{,}680\) candidates.
\item
  Canonical candidate-list SHA-256:
  \texttt{\small 37ee945afb33e5c8\allowbreak f4a63b9874338fb3\allowbreak 8d8cab1e8b1d2c11\allowbreak 3873a2e632ff92f2}.
\end{itemize}

Within this declared family, the canonical \((A,n,k) = (4/3,21,5/2)\) is
the unique candidate within \(0.5\sigma\) and the unique candidate
within \(0.1\sigma\) of the CODATA 2022 central value of \(G\) {[}1{]}.
It is also rank \(1\) against the NIST 2026 central value {[}10{]}.

The five nearest-neighbor rows against CODATA 2022 {[}1{]}:

\begin{table*}[!htb]
\centering\small
\begin{tabular}{@{}
  >{\raggedright\arraybackslash}p{(\linewidth - 12\tabcolsep) * \real{0.1429}}
  >{\raggedright\arraybackslash}p{(\linewidth - 12\tabcolsep) * \real{0.1429}}
  >{\raggedright\arraybackslash}p{(\linewidth - 12\tabcolsep) * \real{0.1429}}
  >{\raggedright\arraybackslash}p{(\linewidth - 12\tabcolsep) * \real{0.1429}}
  >{\raggedright\arraybackslash}p{(\linewidth - 12\tabcolsep) * \real{0.1429}}
  >{\raggedright\arraybackslash}p{(\linewidth - 12\tabcolsep) * \real{0.1429}}
  >{\raggedright\arraybackslash}p{(\linewidth - 12\tabcolsep) * \real{0.1429}}@{}}
\toprule
\begin{minipage}[b]{\linewidth}\raggedright
Rank
\end{minipage} & \begin{minipage}[b]{\linewidth}\raggedright
\(A\)
\end{minipage} & \begin{minipage}[b]{\linewidth}\raggedright
\(n\)
\end{minipage} & \begin{minipage}[b]{\linewidth}\raggedright
\(k\)
\end{minipage} & \begin{minipage}[b]{\linewidth}\raggedright
\(G_{\rm pred}\,\times\!10^{11}\)
\end{minipage} & \begin{minipage}[b]{\linewidth}\raggedright
Residual (ppm)
\end{minipage} & \begin{minipage}[b]{\linewidth}\raggedright
\(z\)-score
\end{minipage} \\
\midrule
\(1\) & \(4/3\) & \(21\) & \(5/2\) & \(6.6743123\) & \(+1.84\) &
\(+0.082\) \\
\(2\) & \(4/3\) & \(21\) & \(3\) & \(6.6500043\) & \(-3640.19\) &
\(-161.97\) \\
\(3\) & \(4/3\) & \(21\) & \(2\) & \(6.6987092\) & \(+3657.19\) &
\(+162.73\) \\
\(4\) & \(4/3\) & \(21\) & \(7/2\) & \(6.6257847\) & \(-7268.96\) &
\(-323.44\) \\
\(5\) & \(4/3\) & \(21\) & \(3/2\) & \(6.7231952\) & \(+7325.90\) &
\(+325.97\) \\
\bottomrule
\end{tabular}
\end{table*}

The same five rows against NIST 2026 {[}10{]} preserve rank order;
canonical residual \(+66.3\) ppm, rank-\(2\) residual \(-3576\) ppm.
Arithmetic rank-loss midpoints between rank \(1\) and the immediate
\(k\)-neighbors \((4/3,21,3)\) and \((4/3,21,2)\) sit at \(-1821.01\)
ppm and \(+1827.67\) ppm relative to \(G_{\rm pred}\). The \(66.3\) ppm
NIST offset is two orders of magnitude smaller than that rank-loss
scale.

Replacing \(m_e\) by \(m_p\) in the same scan produces no candidate
within either \(0.5\sigma\) or \(0.1\sigma\) of the CODATA 2022 central
value.

The reduced-rational prior with \(p\!\le\!5\), \(q\!\le\!30\),
\(p/q\!\in\![1/3,5]\) is the publication declaration. Reproduction
one-liner:

\begin{quote}
\small\ttfamily python3 scan.py\allowbreak\ -{}-benchmark codata2022\allowbreak\ -{}-limit 5
\end{quote}

bundled in the public proof package; expected output reproduces the
candidate-list SHA-256 above and the top-five table.

The bounded scan certifies uniqueness within the declared functional
class with \(m_e\) fixed. It does not certify uniqueness against
arbitrary enlargements of the class; the look-elsewhere problem outside
the declared boundary is treated by the Julia bounded-equation-search
engine of Section 5 (audit row AV6).

\section{3. Algebraic Scaffold}\label{algebraic-scaffold}

\begin{table*}[!htb]
\centering\small
\begin{tabular}{@{}
  >{\raggedright\arraybackslash}p{(\linewidth - 4\tabcolsep) * \real{0.3333}}
  >{\raggedright\arraybackslash}p{(\linewidth - 4\tabcolsep) * \real{0.3333}}
  >{\raggedright\arraybackslash}p{(\linewidth - 4\tabcolsep) * \real{0.3333}}@{}}
\toprule
\begin{minipage}[b]{\linewidth}\raggedright
Coefficient
\end{minipage} & \begin{minipage}[b]{\linewidth}\raggedright
Identification
\end{minipage} & \begin{minipage}[b]{\linewidth}\raggedright
Source
\end{minipage} \\
\midrule
\(21\) & \(\dim(\mathrm{adj}\,SO(7)) = 7\!\cdot\!6/2\) & exact \\
\(5/2\) & \(\rho \cdot e_1 = h^\vee(B_3)/2\), dominant-chamber
convention \(e_1=\omega_1\) & exact, Humphreys {[}3{]} \\
\(4/3\) & \(\dim(SO(8))/\dim(SO(7)) = 28/21\) & exact (primary
anchor) \\
\bottomrule
\end{tabular}
\end{table*}

The \(B_3\) positive root system contains nine roots: three short
positive roots \(e_i\) for \(i=1,2,3\), and six long positive roots
\(e_i\!\pm\!e_j\) for \(1\!\le\!i\!<\!j\!\le\!3\). The Weyl vector
\(\rho = \tfrac{1}{2}\sum_{\alpha>0}\alpha = (5/2, 3/2, 1/2)\) in the
standard orthonormal basis. The basis-invariant statement is
\(\rho\!\cdot\!e_1 = h^\vee(B_3)/2\); the numerical value \(5/2\)
follows in the standard dominant-chamber basis. The dual Coxeter number
is \(h^\vee(B_3) = 5\).

The coefficient \(5/2\) is not \(C_2(\mathrm{vector}, SO(7))\). In the
standard normalization \(C_2(\mathrm{vector}, SO(7)) = 3\); the correct
identification is \(5/2 = h^\vee(B_3)/2 = \rho\!\cdot\!e_1\). Two
further \(B_3\) arithmetic equalities also give \(4/3\) ---
\((h^\vee(B_3)-1)/\mathrm{rank}(B_3) = (5-1)/3 = 4/3\) and
\((\mathrm{rank}(B_3)+1)/\mathrm{rank}(B_3) = (3+1)/3 = 4/3\) --- but
these are not independent invariant derivations; the primary anchor for
\(4/3\) is the \(SO(8)/SO(7)\) embedding ratio.

The scaffold provides exact identifications for the discrete content of
the formula. It does not provide a dynamical derivation. Section 4
documents that no surveyed mechanism class closes the dictionary from
the scaffold to the physical appearance of \(\alpha\), \(m_e\), and the
exponential dressing.

\section{4. Mechanism Gap: Surveyed
Routes}\label{mechanism-gap-surveyed-routes}

The following derivation classes were surveyed and did not produce a
mechanism reproducing the full formula. Each row lists the expected
contribution under that class, the observed surveyed output, and the
structural obstruction.

\begin{table*}[!htb]
\centering\small
\begin{tabular}{@{}
  >{\raggedright\arraybackslash}p{(\linewidth - 6\tabcolsep) * \real{0.2500}}
  >{\raggedright\arraybackslash}p{(\linewidth - 6\tabcolsep) * \real{0.2500}}
  >{\raggedright\arraybackslash}p{(\linewidth - 6\tabcolsep) * \real{0.2500}}
  >{\raggedright\arraybackslash}p{(\linewidth - 6\tabcolsep) * \real{0.2500}}@{}}
\toprule
\begin{minipage}[b]{\linewidth}\raggedright
Route class
\end{minipage} & \begin{minipage}[b]{\linewidth}\raggedright
Expected contribution
\end{minipage} & \begin{minipage}[b]{\linewidth}\raggedright
Surveyed output
\end{minipage} & \begin{minipage}[b]{\linewidth}\raggedright
Obstruction
\end{minipage} \\
\midrule
Heat-kernel / Weyl-character / localization on \(B_3\) & \(\alpha^{21}\)
counting from \(B_3\)-orbit data & \(\alpha^5\exp(-5\alpha/2)\) at
leading order along \(e_1\) & Active-root count along \(e_1\) is \(5\);
\(SO(7)\) adjoint real-orbit dimension ceiling is \(18 < 21\) \\
Instanton / Yang-Mills saddle & Finite exponential dressing
\(\exp(-c\alpha)\) & \(\exp(-c/\alpha)\) & Semiclassical action places
\(\alpha\) in the denominator \\
Soft-photon resummation (YFS / Sudakov / threshold) & Linear-\(\alpha\)
exponentiation & \(\exp(\alpha\cdot C\cdot\log)\), energy-scale logs &
Logarithms in energy scale instead of constant \(\exp(-5\alpha/2)\) \\
Functional determinant on \(S^7\) & Link \(S^7\) determinant to
\(\alpha^{21}\exp(-5\alpha/2)\) & No gauge zero-mode support; wrong
determinant size and sign & \(H^1(S^7) = 0\) removes zero-mode
structure \\
One-loop \(\beta\)-function / threshold coefficient & \(5/2\) from
standard running & \(\beta_0 = 55/3\) (pure YM, \(B_3\));
\(\beta_0 = 15\) (\(N\!=\!1\) SYM, \(B_3\)) & Neither standard
normalization equals \(5/2\) \\
\(SO(7)^\pm\) gauged-supergravity vacua & \(SO(7)\) adjoint structure
from \(N\!=\!8\) supergravity & Both \(SO(7)^\pm\) vacua unstable;
stable \(N\!=\!1\) vacuum is \(G_2\), not \(SO(7)\) & de Wit / Nicolai
{[}4{]}; Borghese et al.~{[}5{]}; Fischbacher {[}6{]} \\
Induced gravity / vacuum response & Electron-mass normalization from
vacuum response & Sakharov-class arguments do not single out \(m_e\) &
Visser {[}7{]}; Chaichian et al.~{[}8{]} \\
Kalinski QED-vacuum kernel & Electron-scale normalization & Bare kernel
\((4/3)(\hbar c/m_e^2)\alpha^{21}\) with fitted \(2^{-1/38}\) factor &
No \(B_3/SO(7)\)-anchored derivation of \(\exp(-5\alpha/2)\) {[}9{]} \\
\bottomrule
\end{tabular}
\end{table*}

The three irreducible anchors a future microscopic theory would need to
close simultaneously:

\begin{enumerate}
\def\labelenumi{\arabic{enumi}.}
\tightlist
\item
  \textbf{\(\alpha^{21}\) counting.} A physical state space or dynamical
  system that produces \(21\) independent modes mapped onto the
  \(SO(7)\) adjoint, rather than analogizing after the fact. The
  orbit-dimension ceiling \(18 < 21\) rules out localization on \(B_3\)
  along any root ray.
\item
  \textbf{\(\exp(-5\alpha/2)\) damping.} A dynamical process producing a
  constant exponential factor with \(\alpha\) in the numerator at
  \(\alpha = \alpha_{\rm EM}\). Surveyed semiclassical actions place
  \(\alpha\) in the denominator; perturbative expansions give
  polynomials in \(\alpha\); soft resummation gives logarithms in energy
  scale.
\item
  \textbf{Electron-mass normalization.} A reason why \(\hbar c/m_e^2\)
  is the gravitational reference scale. The \(m_p\)-substituted bounded
  scan produces no comparable hit. The same locked \(B_3/SO(7)\)
  algebraic scaffold also fails a pre-registered lepton mass-ratio
  holdout: no parameter-free candidate reproduces both \(m_\mu/m_e\) and
  \(m_\tau/m_e\), and the best parameter-free muon-ratio near-hit misses
  by \(1408\) ppm (\(\sim\!760\times\) worse than the \(G\) residual).
\end{enumerate}

The exclusion is over surveyed standard classes, not over all logically
possible theories.

\section{5. Adversarial Validation
Audit}\label{adversarial-validation-audit}

The relation was attacked along seven independent stress modes. Each row
is conditional on the declared bounded prior and engine configuration
cited.

\begin{table*}[!htb]
\centering\small
\begin{tabular}{@{}
  >{\raggedright\arraybackslash}p{(\linewidth - 6\tabcolsep) * \real{0.2500}}
  >{\raggedright\arraybackslash}p{(\linewidth - 6\tabcolsep) * \real{0.2500}}
  >{\raggedright\arraybackslash}p{(\linewidth - 6\tabcolsep) * \real{0.2500}}
  >{\raggedright\arraybackslash}p{(\linewidth - 6\tabcolsep) * \real{0.2500}}@{}}
\toprule
\begin{minipage}[b]{\linewidth}\raggedright
ID
\end{minipage} & \begin{minipage}[b]{\linewidth}\raggedright
Stress mode
\end{minipage} & \begin{minipage}[b]{\linewidth}\raggedright
Declared scope
\end{minipage} & \begin{minipage}[b]{\linewidth}\raggedright
Result
\end{minipage} \\
\midrule
AV0 & Frozen prior & \(22{,}680\)-candidate prior, declared \(A,n,k\)
menus, candidate-list SHA-256 fixed & Canonical \((4/3,21,5/2)\) rank
\(1\); two independent code paths reproduce list and residual \\
AV1 & Sham-target tribunal & \(N\!=\!10{,}000\) synthetic targets under
AV0 grammar & True target rank \(7/10001\) on residual, \(4/10001\)
under strict-cap comparator; surviving competitors concentrate on the
\(n\!=\!21\) shelf \\
AV2 & Source/metrology ablation & Inputs perturbed within reported
uncertainties inside AV0 family & Canonical row shifts \(\le 0.017\) ppm
/ \(\le 0.0008\sigma\); rank and hit counts preserved \\
AV3 & Scaffold invariant census & Frozen-complexity invariant scoring
across tested single-algebra sets & \(B_3/SO(7)\) canonical scaffold
ranks \(1\) \\
AV4 & Blind scaffold prior audit & Independent blind algebraic scoring
across alternative single-algebra scaffolds & \(B_3/SO(7)\) remains top
single-algebra scaffold; \(4/3\) arithmetic echoes recorded as Section 3
echoes \\
AV5 & Lepton mass-ratio holdout & Pre-registered \(m_\mu/m_e\) and
\(m_\tau/m_e\) under locked \(B_3/SO(7)\) scaffold; \(18\)-framework
survey & Ruled out; best parameter-free muon-ratio near-hit \(1408\)
ppm; Koide {[}2{]} charged-lepton relation (\(2/3\)) remains descriptor
only \\
AV6 & Bounded/frontier symbolic null & Julia bounded-equation-search
engine with canonical \(\mathrm{alpha\_family}\) pre-seeded; cost shells
\(C\!=\!18,19,20\); evaluation budgets \(\{20, 30, 40\}\);
\(256\)-target sham-parity & Rank-\(1\) survival under adversarial
competition: among \(2.02\!\times\!10^9\) raw expressions
(\(2.46\!\times\!10^6\) canonical-unique) at \(C\!\le\!20\), the
canonical formula is rank \(1\) at budgets \(30\) and \(40\) with
\(\mathrm{rel\_err}=1.85\!\times\!10^{-6}\) invariant across
\(C\!=\!18,19,20\) (canonical internal engine cost \(28\), not evaluable
at budget \(20\); budget-\(20\) leader is a non-\(\alpha\) pow at
\(\mathrm{rel\_err}\!\sim\!10^{-4}\)). Sham parity \(0/256\) at
\(C\!=\!18\), \(1/256\) at \(C\!=\!19,20\) (non-\(\alpha\) TRASH at
sham-index \(28\)). \\
AV7 & CODATA/NIST target stability & Same frozen prior re-evaluated
against CODATA 2022 {[}1{]} and NIST 2026 {[}10{]} central values & Rank
\(1\) against both targets; rank-loss midpoints \(-1821\) ppm and
\(+1828\) ppm; \(\sim\!13.25\) ppm \(5\sigma\) falsification width \\
\bottomrule
\end{tabular}
\end{table*}

\subsubsection{AV6 Engine Compute Trail}\label{av6-engine-compute-trail}

The Julia bounded-equation-search engine
(\texttt{run\_av6\_exhaustive\_c40\_jl15.jl}, source in the public proof
package) implements grouped/cross-merge canonical reduction, weighted
dynamic chunking, and per-target sham-parity bookkeeping. Canonical
reduction performs in-enumeration equality saturation: expressions are
quotiented modulo associativity and commutativity of \(+,\,\cdot\,\),
multiplicative inverse identities, \(\alpha\)-degree normalization, and
exponent collection, so the \(2.46\!\times\!10^6\) canonical-unique
count at \(C\!=\!20\) is a distinct-expression count, not raw chunk
output. Idempotence of the canonical form and hash stability across
re-runs are checked at engine startup; the \(256\)-target sham-parity
probe doubles as an adversarial engine-level test against accidental
matches. Per-cost shell statistics at completion:

\begin{table*}[!htb]
\centering\small
\begin{tabular}{@{}
  >{\raggedright\arraybackslash}p{(\linewidth - 8\tabcolsep) * \real{0.2000}}
  >{\raggedright\arraybackslash}p{(\linewidth - 8\tabcolsep) * \real{0.2000}}
  >{\raggedright\arraybackslash}p{(\linewidth - 8\tabcolsep) * \real{0.2000}}
  >{\raggedright\arraybackslash}p{(\linewidth - 8\tabcolsep) * \real{0.2000}}
  >{\raggedright\arraybackslash}p{(\linewidth - 8\tabcolsep) * \real{0.2000}}@{}}
\toprule
\begin{minipage}[b]{\linewidth}\raggedright
Shell completion \(C\)
\end{minipage} & \begin{minipage}[b]{\linewidth}\raggedright
Wall time (24 threads)
\end{minipage} & \begin{minipage}[b]{\linewidth}\raggedright
Cumulative raw
\end{minipage} & \begin{minipage}[b]{\linewidth}\raggedright
Cumulative canonical-unique
\end{minipage} & \begin{minipage}[b]{\linewidth}\raggedright
\(\mathrm{sham\_match\_or\_beat}\) at budget \(40\)
\end{minipage} \\
\midrule
\(18\) (publication baseline, 2026-05-08) & \(\sim\!4.5\) h &
\(3.11\!\times\!10^8\) & \(1.30\!\times\!10^6\) & \(0/256\) \\
\(19\) (extending evidence, 2026-05-09) & \(\sim\!14\) h &
\(8.17\!\times\!10^8\) & \(1.80\!\times\!10^6\) & \(1/256\) \\
\(20\) (extending evidence, 2026-05-10) & \(\sim\!28\) h &
\(2.02\!\times\!10^9\) & \(2.46\!\times\!10^6\) & \(1/256\) \\
\bottomrule
\end{tabular}
\end{table*}

With the canonical \(\mathrm{alpha\_family}\) pre-seeded into the engine
grammar, no competitor among the \(2.02\!\times\!10^9\) raw expressions
(\(2.46\!\times\!10^6\) canonical-unique after grammar reduction) at
\(C\!\le\!20\) matches or beats the canonical
\((4/3)\alpha^{21}\exp(-5\alpha/2)\) at evaluation budgets \(30\) and
\(40\). The \(256\)-target sham-parity probe at the full \(C\!\le\!20\)
enumeration returns \(1/256\) --- a single non-\(\alpha\) TRASH pow
expression --- for a per-target hit rate \(\sim\!4\!\times\!10^{-3}\) on
the observed CODATA residual under the same grammar and cost class; the
\(C\!=\!18\) publication baseline returns \(0/256\) at the same
threshold.

The single sham drift to \(1/256\) at \(C \ge 19\) is the same
\texttt{pow} expression at \texttt{sham\_index\ =\ 28}:
\((\exp((5)-((4)^2)\!\cdot\!\pi))^3\) at \(C\!=\!19\) and
\((\exp((1)-(3\!\cdot\!4)\!\cdot\!\pi))^3\) at \(C\!=\!20\), both
classified as \(\mathrm{TRASH\_BROAD\_GRAMMAR}\) and carrying no
\(\alpha\) dependence.

Budget-\(20\) evaluations do not place the canonical at rank \(1\)
because its internal engine cost is \(28\) and exceeds the budget cap;
the canonical is not assigned a target-fit score under that constraint.
The budget-\(20\) leader is a non-\(\alpha\) pow expression with
\(\mathrm{rel\_err}\!\sim\!4.6\!\times\!10^{-4}\) at \(C\!=\!18\) and
\(\sim\!1.06\!\times\!10^{-4}\) at \(C\!=\!19,20\). At budgets \(30\)
and \(40\) the canonical evaluates and dominates by three orders of
magnitude in \(\mathrm{rel\_err}\).

\(C\!=\!18\) is the publication baseline; \(C\!=\!19\) and \(C\!=\!20\)
extend sibling coverage and sham-parity robustness while preserving the
same canonical rank-\(1\) outcome and \(\mathrm{rel\_err}\).

The \(\mathrm{jl15}\) engine supersedes two earlier internal runners.
The intermediate \(\mathrm{jl14}\) had an \(\mathrm{alpha\_family}\)
seed envelope that excluded the cost-\(28\) canonical from the generated
set; the omission was caught during pre-publication ledger audit and
corrected in \(\mathrm{jl15}\) (seed envelope raised to
\(\max(\mathrm{budgets})\)). The earlier \(\mathrm{jl13}\) used a
coarser sibling-class classifier and is not part of the reported
ledgers. Both prior runners are present in the engine source tree for
inspection; their outputs are not aggregated with the \(C\!=\!18,19,20\)
\(\mathrm{jl15}\) ledgers reported here. Detection of the
\(\mathrm{jl14}\) omission during audit is itself a positive datum on
engine adversariality: the canonical was not granted to the engine by
construction; it had to survive a generation-side seeding bug that
initially hid it.

Three further audit paths were considered but not executed because they
do not change the verdict at first order. (i) Extending shell completion
beyond \(C\!=\!20\) roughly doubles wall time per shell
(\(C\!=\!18\!\sim\!4.5\) h, \(C\!=\!19\!\sim\!14\) h,
\(C\!=\!20\!\sim\!28\) h on 24 threads); given canonical rank \(1\) with
identical \(\mathrm{rel\_err}=1.85\!\times\!10^{-6}\) is already
invariant across \(C\!=\!18,19,20\), additional shells would extend
sibling coverage and sham-parity robustness rather than flip the
outcome. (ii) Broadening the grammar with non-\(\alpha\) primitives
enlarges the sham pool but cannot remove a pre-seeded canonical from
rank \(1\) since comparison is by absolute residual. (iii) Removing the
\(\mathrm{alpha\_family}\) pre-seed tests whether brute-force
enumeration alone reaches a cost-\(28\) canonical within available
compute, a different question than adversarial competition against a
present canonical; the reported \(C\!=\!18,19,20\) enumeration already
answers the cost-coverage question through the explicit budget-\(20\)
versus budget-\(30/40\) contrast.

\section{6. Metrology Residuals}\label{metrology-residuals}

Predicted
\(G_{\rm pred} = 6.6743123 \times 10^{-11}\,\mathrm{m^3\,kg^{-1}\,s^{-2}}\)
compared against modern precision measurements and adjustments:

\begin{table*}[!htb]
\centering\small
\begin{tabular}{@{}
  >{\raggedright\arraybackslash}p{(\linewidth - 14\tabcolsep) * \real{0.1250}}
  >{\raggedright\arraybackslash}p{(\linewidth - 14\tabcolsep) * \real{0.1250}}
  >{\raggedright\arraybackslash}p{(\linewidth - 14\tabcolsep) * \real{0.1250}}
  >{\raggedright\arraybackslash}p{(\linewidth - 14\tabcolsep) * \real{0.1250}}
  >{\raggedright\arraybackslash}p{(\linewidth - 14\tabcolsep) * \real{0.1250}}
  >{\raggedright\arraybackslash}p{(\linewidth - 14\tabcolsep) * \real{0.1250}}
  >{\raggedright\arraybackslash}p{(\linewidth - 14\tabcolsep) * \real{0.1250}}
  >{\raggedright\arraybackslash}p{(\linewidth - 14\tabcolsep) * \real{0.1250}}@{}}
\toprule
\begin{minipage}[b]{\linewidth}\raggedright
Source
\end{minipage} & \begin{minipage}[b]{\linewidth}\raggedright
Citation
\end{minipage} & \begin{minipage}[b]{\linewidth}\raggedright
\(G\,\times\!10^{11}\)
\end{minipage} & \begin{minipage}[b]{\linewidth}\raggedright
\(\sigma\,\times\!10^{11}\)
\end{minipage} & \begin{minipage}[b]{\linewidth}\raggedright
Type
\end{minipage} & \begin{minipage}[b]{\linewidth}\raggedright
Method
\end{minipage} & \begin{minipage}[b]{\linewidth}\raggedright
Residual (ppm)
\end{minipage} & \begin{minipage}[b]{\linewidth}\raggedright
\(z\)-score
\end{minipage} \\
\midrule
CODATA 2014 & Mohr et al.~(2016) & \(6.67408\) & \(0.00031\) &
Adjustment & LSQ & \(+34.80\) & \(+0.749\) \\
CODATA 2018 & Tiesinga et al.~(2021) & \(6.67430\) & \(0.00015\) &
Adjustment & LSQ & \(+1.84\) & \(+0.082\) \\
CODATA 2022 & Mohr et al.~(2025) {[}1{]} & \(6.67430\) & \(0.00015\) &
Adjustment & LSQ & \(+1.84\) & \(+0.082\) \\
UWash & Gundlach \& Merkowitz (2000) & \(6.674215\) & \(0.000092\) &
Primary & Torsion balance, AAF & \(+14.58\) & \(+1.057\) \\
BIPM 2001 & Quinn et al.~(2001) & \(6.67559\) & \(0.00027\) & Primary &
Torsion strip, electrostatic & \(-191.40\) & \(-4.732\) \\
UZur-06 & Schlamminger et al.~(2006) & \(6.67425\) & \(0.00012\) &
Primary & Beam balance & \(+9.33\) & \(+0.519\) \\
JILA-10 & Parks \& Faller (2010) & \(6.67234\) & \(0.00014\) & Primary &
Fabry-Perot & \(+295.59\) & \(+14.088\) \\
BIPM 2013 & Quinn et al.~(2013) & \(6.67554\) & \(0.00016\) & Primary &
Torsion strip + TOS & \(-183.92\) & \(-7.673\) \\
Atom interferometry & Rosi et al.~(2014) & \(6.67191\) & \(0.00099\) &
Primary & Atom interferometry & \(+360.06\) & \(+2.427\) \\
UCI-14 & Newman et al.~(2014) & \(6.67433\) & \(0.00013\) & Primary &
Torsion balance, TOS & \(-2.65\) & \(-0.136\) \\
HUST-18 TOS & Li et al.~(2018) & \(6.674184\) & \(0.000078\) & Primary &
Torsion balance, TOS & \(+19.22\) & \(+1.645\) \\
HUST-18 AAF & Li et al.~(2018) & \(6.674484\) & \(0.000078\) & Primary &
Torsion balance, AAF & \(-25.72\) & \(-2.201\) \\
NIST 2026 & Schlamminger et al.~(2026) {[}10{]} & \(6.67387\) &
\(0.00038\) & Primary & BIPM torsion at NIST & \(+66.27\) &
\(+1.164\) \\
\bottomrule
\end{tabular}
\end{table*}

Sign convention: residual \(= (G_{\rm pred} - G)\times 10^6 / G\) in
ppm; \(z = (G_{\rm pred} - G)/\sigma\). Positive entries indicate the
formula prediction lies above the measured central value.

The CODATA adjustments aggregate the primary measurements through
least-squares treatment. The spread among primary determinations (BIPM
2001 / 2013 on the high side, JILA-10 / Rosi on the low side) is visible
in the residual column.

\section{7. Prior Art Boundary}\label{prior-art-boundary}

The public Rochester homepage of M. Kalinski {[}9{]} contains the bare
kernel \((4/3)(\hbar c/m_e^2)\alpha^{21}\) with a numerically fitted
correction factor \(2^{-1/38}\). An archived snapshot of that page and a
bundled local capture are part of the public proof package
(\texttt{public\_proof\_package/\allowbreak kalinski\_rochester\_capture\_2026-04-23.md}).

The present manuscript adds:

\begin{enumerate}
\def\labelenumi{\arabic{enumi}.}
\tightlist
\item
  The \(B_3/SO(7)/SO(8)\) algebraic scaffold for the discrete
  coefficients \(21\), \(5/2\), and \(4/3\).
\item
  The replacement of the fitted \(2^{-1/38}\) factor by the
  \(B_3\)-anchored \(\exp(-5\alpha/2)\), tied to
  \(h^\vee(B_3)/2 = \rho\!\cdot\!e_1\).
\item
  The bounded \(22{,}680\)-candidate reproducibility scan with declared
  prior and candidate-list SHA-256.
\item
  The structured mechanism-gap survey of Section 4.
\item
  The AV0-AV7 adversarial validation audit of Section 5, including the
  Julia bounded-equation-search engine through \(C\!=\!20\).
\end{enumerate}

\section{8. Reproducibility}\label{reproducibility}

The public proof package bundles:

\begin{itemize}
\tightlist
\item
  \texttt{scan.py}, \texttt{canonical\_bounded\_prior.py},
  \texttt{constants.json}, \texttt{benchmarks.json},
  \texttt{declared\_prior.csv}: the bounded \(22{,}680\)-candidate scan,
  reproducible on a laptop in seconds.
\item
  \texttt{top\_candidates.csv}, \texttt{nist2026\_top\_candidates.csv}:
  the five nearest-neighbor tables of Section 2.
\item
  \texttt{av6\_engine/\allowbreak run\_av6\_\allowbreak exhaustive\_c40\_jl15.jl},
  \texttt{av6\_engine/\allowbreak launch\_jl15.sh},
  \texttt{av6\_engine/\allowbreak spec.json}:
  the Julia bounded-equation-search engine source and launcher.
\item
  \texttt{av6\_outputs/c18/}, \texttt{av6\_outputs/c19/},
  \texttt{av6\_outputs/c20/}: per-cost ledgers
  \texttt{bounded\_shell\_stats.csv},
  \texttt{same\_target\_rank\_tables.csv},
  \texttt{sham\_best\_ledger.csv} directly auditable without re-running
  the engine. The \(C\!=\!18\) directory additionally contains the
  publication-grade KPI verdict file
  \texttt{diff\_ledgers\_publication\_verdict.md}.
\item
  \texttt{metrology\_residuals.csv},
  \texttt{mechanism\_gap\_routes.csv}: tabular sources of Sections 6 and
  4.
\item
  \texttt{manifest.json}: SHA-256 hashes for every package file plus the
  canonical candidate-list hash.
\item
  \texttt{LICENSE}, \texttt{SCOPE.md}, \texttt{README.md}: reuse rights
  and explicit scope statement.
\end{itemize}

Three independent verification paths:

\begin{enumerate}
\def\labelenumi{\arabic{enumi}.}
\tightlist
\item
  \textbf{Numerical.} Recompute \(G_{\rm pred}\) from CODATA 2022
  {[}1{]} values of \(\alpha\), \(m_e\), \(\hbar\), \(c\), then verify
  the quoted residual and \(z\)-score.
\item
  \textbf{Algebraic.} Verify \(\dim(\mathrm{adj}\,SO(7)) = 21\),
  \(h^\vee(B_3) = 5\), \(\rho\!\cdot\!e_1 = 5/2\), and
  \(\dim(SO(8))/\dim(SO(7)) = 28/21 = 4/3\) from \(B_3\) root data
  {[}3{]}.
\item
  \textbf{Bounded-statistical.} Run
  \texttt{python3\ scan.py\ -\/-benchmark\ codata2022\ -\/-limit\ 5};
  expected output reproduces candidate-list SHA-256
  \texttt{37ee945afb33e5c8f4a63b9874338fb3\allowbreak 8d8cab1e8b1d2c113873a2e632ff92f2}
  and the top-five table of Section 2. For the AV6 engine,
  \texttt{./av6\_engine/launch\_jl15.sh\ 18} reproduces the publication
  baseline in \(\sim\!4.5\) hours on \(24\) threads.
\end{enumerate}

\section{Statement of AI Authorship}\label{ai-authorship}

\noindent\textbf{This research and the resulting manuscript were carried out by an orchestration of AI agents under the named author's direction.} The author is an independent researcher without specialist training in theoretical physics. His role was vision, problem framing, target selection, and gate-keeping decisions: which leads to pursue, which to drop, when to escalate to multi-model adversarial review, and when a draft was acceptable. He did not personally derive the relation, identify the algebraic scaffold, design or operate the bounded-equation-search engine, or write the prose of this paper.

\medskip

The technical work --- empirical fit of \(G\) to the \(\alpha^{21}\exp(-5\alpha/2)\) family; identification of the \(B_3/SO(7)/SO(8)\) algebraic identities; design and execution of the bounded-prior reproducibility scan; implementation and runtime of the Julia bounded-equation-search engine and its sham-parity controls; structured survey of candidate mechanism routes; assembly of the public proof package; and the writing of this manuscript --- was performed by an orchestration of OpenAI GPT-5.4 / GPT-5.5, Google Gemini 3 Pro, Anthropic Claude Opus 4.7 and Sonnet 4.6, and Deepseek V4. Multi-model adversarial review served as the primary quality gate; council disagreements were surfaced to the named author for adjudication.

\medskip

The named author chose the research target, accepted or rejected lines of work, approved every draft revision, and bears scientific responsibility for the content, claims, and citations herein.

\section{References}\label{references}

\begin{sloppypar}
\begin{enumerate}
\def\labelenumi{\arabic{enumi}.}
\item
  P. J. Mohr, D. B. Newell, B. N. Taylor, and E. Tiesinga, ``CODATA
  recommended values of the fundamental physical constants: 2022,''
  \emph{Rev.~Mod. Phys.} \textbf{97}, 025002 (2025). DOI:
  \url{https://doi.org/10.1103/RevModPhys.97.025002}.
\item
  Y. Koide, ``New view of quark and lepton mass hierarchy,'' \emph{Phys.
  Rev.~D} \textbf{28}, 252--254 (1983). DOI:
  \url{https://doi.org/10.1103/PhysRevD.28.252}.
\item
  J. E. Humphreys, \emph{Introduction to Lie Algebras and Representation
  Theory}, Graduate Texts in Mathematics Vol. 9 (Springer, New York,
  1972). DOI: \url{https://doi.org/10.1007/978-1-4612-6398-2}.
\item
  B. de Wit and H. Nicolai, ``\(N=8\) supergravity with local
  \(SO(8)\times SU(8)\) invariance,'' \emph{Phys. Lett. B} \textbf{108},
  285--290 (1982). DOI: \url{https://doi.org/10.1016/0370-2693(82)91194-7}.
\item
  A. Borghese, A. Guarino, and D. Roest, ``Triality, Periodicity and
  Stability of \(SO(8)\) Gauged Supergravity,'' \emph{JHEP} \textbf{05}
  (2013) 107. DOI: \url{https://doi.org/10.1007/JHEP05(2013)107}; arXiv:
  \url{https://arxiv.org/abs/1302.6057}.
\item
  T. Fischbacher, ``The Encyclopedic Reference of Critical Points for
  \(SO(8)\)-Gauged \(N=8\) Supergravity Part 1: Cosmological Constants
  in the Range \(-\Lambda/g^2 \in [6;14.7)\),'' arXiv:
  \url{https://arxiv.org/abs/1109.1424} [hep-th].
\item
  M. Visser, ``Sakharov's induced gravity: a modern perspective,''
  \emph{Mod. Phys. Lett. A} \textbf{17}, 977--992 (2002). DOI:
  \url{https://doi.org/10.1142/S0217732302006886}; arXiv:
  \url{https://arxiv.org/abs/gr-qc/0204062}.
\item
  M. Chaichian, M. Oksanen, and A. Tureanu, ``Sakharov's induced gravity
  and the Poincaré gauge theory,'' in \emph{Jacob Bekenstein: The
  Conservative Revolutionary}, ed.~L. Brink, V. Mukhanov, E. Rabinovici,
  and K. K. Phua (World Scientific, Singapore, 2019), pp.~271--299. DOI:
  \url{https://doi.org/10.1142/9789811203961_0021}; arXiv:
  \url{https://arxiv.org/abs/1805.03148}.
\item
  M. Kalinski, ``Matt Kalinski's homepage,'' University of Rochester,
  gravitational-constant note beginning ``Relation between the
  Gravitational constant and the Planck constant,'' accessed April 23,
  2026; archived snapshot:
  \url{https://web.archive.org/web/20250408212219/https://www.pas.rochester.edu/~mkal/};
  bundled local capture:
  \texttt{public\_proof\_package/\allowbreak kalinski\_rochester\_capture\_2026-04-23.md}.
\item
  National Institute of Standards and Technology, ``NIST Weighs In on
  the Mystery of the Gravitational Constant,'' April 16, 2026. URL:
  \url{https://www.nist.gov/news-events/news/2026/04/nist-weighs-mystery-gravitational-constant};
  S. Schlamminger, L. Chao, V. Lee, C. Shakarji, A. Possolo, D. Newell,
  J. Stirling, R. Cochran, and C. Speake, ``Redetermination of the
  gravitational constant with the BIPM torsion balance at NIST,''
  \emph{Metrologia} (published online April 16, 2026). DOI:
  \url{https://doi.org/10.1088/1681-7575/ae570f}.
\end{enumerate}
\end{sloppypar}

\end{document}
